\section*{1. Constrained legal-action control}

Let $I_t=\Pi_{\mathrm{act}}(o_{\leq t})$ be everything visible to the acting
seat: its own hand and private state, the public board and history, and the
current simulator menu. Opponent hand identities, either deck order,
unrevealed Prize identities, terminal results not yet observed, and future
transitions are excluded. The pinned simulator, not a learned rule model,
defines
\[
  \mathcal A_t=\mathcal A(I_t).
\]

A prompt with $n$ options and legal selection size $\ell\leq k\leq u$
contains
\[
  |\mathcal A_t|
  =\sum_{k=\ell}^{u} P(n,k)
  =\sum_{k=\ell}^{u}\frac{n!}{(n-k)!}
\]
ordered complete actions. Enumerating that set is unnecessary. From prefix
$p$, the decoder considers every unused extension $p\mathbin{\Vert}i$ and,
once $|p|\geq\ell$, an explicit \textsc{stop}. At $|p|=u$, the prefix is
complete. Thus, for $a=(i_1,\ldots,i_k)$,
\[
 \pi_\theta(a\mid I_t)
 =\prod_{j=1}^{k}\pi_\theta(i_j\mid I_t,p_{j-1})
  \;\pi_\theta(\textsc{stop}\mid I_t,p_k)^{\mathbf 1[k<u]} .
\]
Every stage is masked to its simulator-supplied candidate set $C_t(p)$:
\[
 \pi_\theta(i\mid I_t,p)
 =\frac{\mathbf 1[i\in C_t(p)]e^{\ell_i}}
        {\sum_{j\in C_t(p)}e^{\ell_j}} .
\]
This factorization preserves complete legal support and recomputes legality
after each selection, including order-sensitive effects.
